\exercisesection{Composition Product}

\begin{enumerate}
\def\labelenumi{\arabic{enumi})}
\tightlist
\item
  Let \(k\) be a commutative ring, A an associative unital \(k\) -algebra, M an (A, A)-bimodule. An extension of A by M is the data of an associative unital \(k\) -algebra B and of an exact sequence of \(k\) -modules
\end{enumerate}

\[
(\mathcal {E}): 0 \rightarrow \mathrm{M} \stackrel {{i}} {{\rightarrow}} \mathrm{B} \stackrel {{p}} {{\rightarrow}} \mathrm{A} \rightarrow 0,
\]

where p is an algebra homomorphism and where

\[
i (p (b) m) = b i (m) \quad \text { and } \quad i (m p (b)) = i (m) b \quad \text { for } \quad b \in \mathbf {B}, \quad m \in \mathbf {M}.
\]

Two extensions \(0 \to M \xrightarrow{i} B \xrightarrow{p} A \to 0\) and \(0 \to M \xrightarrow{i'} B' \xrightarrow{p'} A \to 0\) are said to be equivalent if there exists an algebra homomorphism \(u: B \to B'\) such that \(p' \circ u = p\) and \(u \circ i = i'\) ; the homomorphism \(u\) is then bijective.

\begin{enumerate}
\def\labelenumi{\alph{enumi})}
\item
  We assume that the exact sequence of \(k\) -modules \(0 \to \mathbf{M} \xrightarrow{i} \mathbf{B} \xrightarrow{p} \mathbf{A} \to 0\) is split; let \(s: \mathbf{A} \to \mathbf{B}\) be a \(k\) -linear section of \(p\) . Show that there exists an element \(f \in \mathbf{C}^2(\mathbf{A}, \mathbf{M})\) (X, p.~195, exercise 26, b)) such that \(s(a) s(a') - s(aa') = i(f(a, a'))\) for \(a, a' \in \mathbf{A}\) . Show that \(f\) is a 2-cocycle and that its class in \(\mathbf{H}^2(\mathbf{A}, \mathbf{M})\) is independent of the choice of \(s\) ; we denote it \(c(\mathcal{E})\) .
\item
  Let \(\mathrm{Ex}_{k}(\mathbf{A},\mathbf{M})\) be the set of equivalence classes of extensions of A by M which are trivial as extensions of k-modules. Show that c defines a bijection from \(\mathrm{Ex}_{k}(\mathbf{A},\mathbf{M})\) onto \(\mathrm{H}^{2}(\mathbf{A},\mathbf{M})\) .
\item
  Let \(\mathcal{E}:0\to\mathbf{M}\stackrel{i}{\to}\mathbf{B}\stackrel{p}{\to}\mathbf{A}\to0\) and \(\mathcal{E}':0\to\mathbf{M}\stackrel{i'}{\to}\mathbf{B}\stackrel{p'}{\to}\mathbf{A}\to0\) be two extensions of A by M. We denote by C the subalgebra of \(\mathbf{B}\times\mathbf{B}'\) consisting of the elements \((b,b')\) such that \(p(b)=p'(b')\) , c the ideal of C consisting of the elements \((i(m),-i'(m))\) for \(m\in\mathbf{M}\) , \(\mathbf{B}''\) the algebra C/c, \(\pi:\mathbf{C}\to\mathbf{B}''\) the quotient homomorphism. Let \(i'':\mathbf{M}\to\mathbf{B}''\) and \(p'':\mathbf{B}''\to\mathbf{A}\) be the homomorphisms such that \(i''(m)=\pi((m,0))\) for \(m\in\mathbf{M}\) and \(p''\pi((b,b'))=p(b)\) for \(b\in\mathbf{B},b'\in\mathbf{B}'\) . Show that the sequence \(0\to\mathbf{M}\to\mathbf{B}''\to\mathbf{A}\to0\) is an extension of A by M, called the sum extension of ( \(\mathcal{E}\) ) and ( \(\mathcal{E}'\) ), and that this defines a commutative group law on the set of equivalence classes of extensions of A by M. In particular, the set \(\mathrm{Ex}_{k}(\mathbf{A},\mathbf{M})\) is thereby endowed with a group structure, and the map \(c\) is a group isomorphism.
\end{enumerate}

\begin{enumerate}
\def\labelenumi{\arabic{enumi})}
\setcounter{enumi}{1}
\item
  Suppose that the \(k\) algebra A admits an augmentation \(\pi:\mathbf{A}\to k\) (X, p.~196, exercise 27) and that
\item
  Suppose that the \(k\) -algebra \(\mathbf{A}\) admits an augmentation \(\pi : \mathbf{A} \to k\) (X, p.~196, exercise 27) and that, moreover, \(\mathbf{A}\) is a projective \(k\) -module. We denote by I the kernel of \(\pi\) . Let \(\mathbf{M}\) be a left \(\mathbf{A}\) -module.
\end{enumerate}

\begin{enumerate}
\def\labelenumi{\alph{enumi})}
\item
  Let \((\mathcal{E}):0\to M_{(\pi)}\stackrel{i}{\to}\mathbf{B}\stackrel{p}{\to}\mathbf{A}\to 0\) be an extension of \(\mathbf{A}\) by the \((\mathbf{A},\mathbf{A})\) -bimodule \(M_{(\pi)}\) (X, p.~196, exercise 27); we denote by \(\overline{\mathbf{B}}\) the ideal of \(\mathbf{B}\) of elements \(b\in \mathbf{B}\) such that \(p(b)\in I\) . We have \(i(m)b = 0\) for \(m\in \mathbf{M}\) and \(b\in \overline{\mathbf{B}}\) , so that \(\overline{\mathbf{B}}\) is equipped with a structure of left \(\mathbf{A}\) -module such that \(p(\beta)b = \beta b\) for \(\beta \in \mathbf{B}\) , \(b\in \overline{\mathbf{B}}\) . Show that there exists an exact sequence of \(\mathbf{A}\) -modules \(0\to \mathrm{M}\to \overline{\mathrm{B}}\to I\to 0\) ; we denote by \(\theta (\mathcal{E})\) the class of this exact sequence in \(\operatorname {Ext}_{\mathbf{A}}^{1}(I,\mathbf{M})\) .
\item
  Show that \(\theta\) defines a map from \(\mathrm{Ex}_{k}(\mathbf{A},\mathbf{M}_{(\pi)})\) to \(\mathrm{Ext}_{\mathbf{A}}^{1}(I,\mathbf{M})\) . Prove the commutativity of the diagram:
\end{enumerate}

\[
\begin{array}{c c c} \operatorname{Ex} _ {k} (\mathbf {A}, \mathbf {M} _ {(\pi)}) & \stackrel {{\theta}} {{\longrightarrow}} & \operatorname{Ext} _ {\mathbf {A}} ^ {1} (I, \mathbf {M}) \\ \downarrow^ {c} & & \downarrow^ {- \delta} \\ \mathbf {H} ^ {2} (\mathbf {A}, \mathbf {M} _ {(\pi)}) & \stackrel {{\varphi}} {{\longrightarrow}} & \operatorname{Ext} _ {\mathbf {A}} ^ {2} (k, \mathbf {M}) \end{array}
\]

where δ is the connecting homomorphism deduced from the exact sequence \(0 \rightarrow I \rightarrow A \xrightarrow{\pi} k \rightarrow 0\) , and where φ is the isomorphism obtained by computing with the standard resolution (X, p.~196, exercise 27, b)). Deduce that θ is an isomorphism of groups.

\begin{enumerate}
\def\labelenumi{\arabic{enumi})}
\setcounter{enumi}{2}
\tightlist
\item
  Let G be a group, A the algebra \(\mathbf{Z}^{(\mathrm{G})}, j: \mathrm{G} \to \mathrm{A}\) the canonical injection. One endows A with the augmentation \(\pi: A \to Z\) defined by \(\pi(e_g) = 1\) for all \(g \in G\) . Let M be a left A-module.
\end{enumerate}

\begin{enumerate}
\def\labelenumi{\alph{enumi})}
\item
  Let \(0 \to M_{(\pi)} \xrightarrow{i} B \xrightarrow{p} A \to 0\) be an extension of \(A\) by the \((A, A)\) -bimodule \(M_{(\pi)}\) (X, p.~196, exercise 27). We denote by \(\Gamma\) the subset of \(B\) consisting of the elements \(b\) of \(B\) such that \(p(b) \in j(G)\) . Show that \(\Gamma\) , endowed with the multiplication induced by that of \(B\) , is a group, an extension of \(G\) by the \(G\) -module \(M\) .
\item
  The preceding construction defines a map from \(\mathrm{Ex}_{\mathbf{Z}}(\mathbf{A},\mathbf{M}_{(\pi)})\) to \(\mathrm{Ex}(\mathbf{G},\mathbf{M})\) (VIII, § 13). Show that the diagram
\end{enumerate}

\[
\begin{array}{c} \operatorname{Ex} _ {\mathbf {Z}} (\mathbf {A}, \mathbf {M} _ {(\pi)}) \to \operatorname{Ex} (\mathbf {G}, \mathbf {M}) \\ \downarrow^ {c} \qquad \qquad \qquad \qquad \qquad \qquad \qquad \qquad \qquad \qquad \downarrow \\ \mathbf {H} ^ {2} (\mathbf {A}, \mathbf {M} _ {(\pi)}) \xrightarrow {\varphi^ {2}} \mathbf {H} ^ {2} (\mathbf {G}, \mathbf {M}) \end{array}
\]

is commutative.

* 4) Let g be a Lie k-algebra, M a g-module. An extension of g by M is an exact sequence of k-modules \(0 \rightarrow M \xrightarrow{i} h \xrightarrow{p} g \rightarrow 0\) where h is a Lie k-algebra, p a homomorphism of Lie algebras, and where \(i(p(H)m) = [H, i(m)]\) for \(H \in h, m \in M\) . Two extensions

\[
0 \rightarrow \mathrm{M} \stackrel {i} {\rightarrow} \mathfrak {h} \stackrel {p} {\rightarrow} \mathfrak {g} \rightarrow 0, 0 \rightarrow \mathrm{M} \stackrel {i ^ {\prime}} {\rightarrow} \mathfrak {h} ^ {\prime} \stackrel {p ^ {\prime}} {\rightarrow} \mathfrak {g} \rightarrow 0
\]

are said to be equivalent if there exists a homomorphism \(u: \mathfrak{h} \to \mathfrak{h}'\) such that \(p' \circ u = p\) , \(u \circ i = i'\) . We denote by Liex (g, M) the set of equivalence classes of extensions of g by M. We assume that g is a free k-module.

\begin{enumerate}
\def\labelenumi{\alph{enumi})}
\tightlist
\item
  Let \((\mathcal{E}):0\to \mathrm{M}\stackrel {i}{\to}\mathfrak{h}\stackrel {p}{\to}\mathfrak{g}\to 0\) be an extension of \(\mathfrak{g}\) by \(\mathbf{M}\) ; we choose a \(k\) -linear section \(s:\mathfrak{g}\to \mathfrak{h}\) of \(p\) . We define an element \(f\in C^2 (\mathfrak{g},\mathbf{M})\) (X, p.~196, exercise 28, b)) by setting
\end{enumerate}

\[
i (f (x, y)) = [ s (x), s (y) ] - s ([ x , y ]) \quad \text { for } \quad x  ,   y \in \mathfrak {g}  .
\]

Show that \(f\) is a 2-cocycle and that its class in \(\mathbf{H}^{2}(\mathfrak{g},\mathbf{M})\) does not depend on the choice of \(s\) ; we denote it \(\theta(\mathcal{E})\) . b) Show that \(\theta\) defines a bijection from Liex \((\mathfrak{g},\mathbf{M})\) onto \(\mathbf{H}^{2}(\mathfrak{g},\mathbf{M})\) . Describe directly the group structure on Liex \((\mathfrak{g},\mathbf{M})\) obtained by transport via \(\gamma\) .

\begin{enumerate}
\def\labelenumi{\alph{enumi})}
\setcounter{enumi}{2}
\tightlist
\item
  Let U be the enveloping algebra of g, \(M_{(\pi)}\) the (U, U)-bimodule associated with M by means of the augmentation \(\pi: U \to k\) (X, p.~196, exercise 27). Let \(0 \to M_{(\pi)} \to B \xrightarrow{p} U \to 0\) be an extension of U by \(M_{(\pi)}\) (exercise 1); we denote by b the set of elements b of B such that \(p(b) \in j(g)\) .
\end{enumerate}

Show that b is a Lie k-algebra, an extension of g by M. This construction defines a map \(\lambda: \mathrm{Ex}_{k}(\mathrm{U}, \mathrm{M}_{(\pi)}) \to \mathrm{Liex}(\mathfrak{g}, \mathrm{M})\) . Show that the diagram

\[
\begin{array}{c c c} \operatorname{Ex} _ {k} (\mathrm{U}, \mathrm{M} _ {(\pi)}) & \xrightarrow {\lambda} & \text { Liex(g,M) } \\ \downarrow^ {c} & & \downarrow^ {\gamma} \\ \mathrm{H} ^ {2} (\mathrm{U}, \mathrm{M} _ {(\pi)}) & \xrightarrow {\varphi^ {2}} & \mathrm{H} ^ {2} (\mathrm{g,M}) \end{array}
\]

(where \(\varphi^2\) is the isomorphism defined in X, p.~196, exercise 27) is commutative. Deduce that \(\lambda\) is an isomorphism. *

¶ 5) Let G, F be two groups.

\begin{enumerate}
\def\labelenumi{\alph{enumi})}
\item
  For every extension \(\mathcal{E}: F \to E \to G\) of \(G\) by \(F\) (I, p.~62), the action of \(E\) on \(F\) by inner automorphisms defines a homomorphism \(E \to \text{Aut}(F)\) , whence, by passing to the quotient, a group homomorphism \(\theta: G \to \text{Aut}(F)/\text{Int}(F)\) , called the homomorphism associated with the extension \(\mathcal{E}\) . Two isomorphic extensions of \(G\) by \(F\) have the same associated homomorphism.
\item
  Conversely, we consider a group homomorphism \(\theta: G \to \text{Aut} (F)/\text{Int} (F)\) . Let
\end{enumerate}

\[
p: \operatorname{Aut} (\mathbf {F}) \rightarrow \operatorname{Aut} (\mathbf {F}) / \operatorname{Int} (\mathbf {F})
\]

the quotient map. We choose a map \(\sigma\) from G to Aut (F) such that \(p \circ \sigma = \theta\) , and a map \(f\) from G \(\times\) G to F such that \(\sigma(x)\, \sigma(y)\, (\sigma(xy))^{-1} = \text{Int}(f(x,y))\) for \(x,y \in G\) ; set \(c(x,y,z) = \sigma(x)\, (f(y,z)).f(x,yz).(f(xy,z))^{-1}.(f(x,y))^{-1}\) . Show that \(c(x,y,z)\) belongs to the center Z of F.

\begin{enumerate}
\def\labelenumi{\alph{enumi})}
\setcounter{enumi}{2}
\item
  We consider Z as a G-module by setting \(g . z = \sigma(g) (z)\) for \(g \in G, z \in Z\) ; this definition is independent of the choice of \(\sigma\) . The map c defines an element of \(C^{3}(G, Z)\) ; show that c is a 3-cocycle, and that its class in \(H^{3}(G, Z)\) is independent of the choices of \(\sigma\) and of f. We denote it \(\omega(G, F, \theta)\) . d) Show that the vanishing of \(\omega(G, F, \theta)\) is a necessary and sufficient condition for there to exist an extension of G by F having as associated homomorphism \(\theta\) . If \(\omega(G, F, \theta) = 0\) , show that one can choose f so that the cocycle c is zero, then define a group law on \(G \times F\) .
\item
  Assume \(\omega(G, F, \theta) = 0\) . Show that the set of isomorphism classes of extensions of \(G\) by \(F\) of the associated homomorphism \(\theta\) is a principal homogeneous set under the group H \(^{2}\) (G, Z).
\item
  Let H be a group, M a \(\mathbf{Z}^{(H)}\) -module, \(x\) be an element of \(\mathrm{H}^3 (\mathrm{H},\mathrm{M})\) . Show that there exists a group E with center M and a homomorphism \(\varphi :\mathrm{H}\to \mathrm{Aut}(\mathrm{E}) / \mathrm{Int}(\mathrm{E})\) , inducing on M the given H-module structure, such that \(\omega (\mathrm{H},\mathrm{E},\varphi) = x\) (take \(\mathbf{E} = \mathbf{M}\times \mathbf{L}\) , where L is the free group on \(\mathrm{H}\times \mathrm{H}\) ).
\end{enumerate}

\begin{enumerate}
\def\labelenumi{\arabic{enumi})}
\setcounter{enumi}{5}
\tightlist
\item
  Let M be a right A-module and N a left A-module. For \(n \geqslant 0\) , we denote by \(\mathcal{T}_{n}(M, N)\) the set of triples \((P, p, q)\) , where P is a complex of finitely generated free A-modules, such that \(P_{r} = 0\) for \(r < 0\) and \(r > n\) , and where \(p: P \to N\) and \(q: P^{*} \to M(n)\) are morphisms of A-complexes (we denote by \(P^{*}\) the Homgr complex \(_{A}\) \((P, A)\) ). Two elements \((P, p, q)\) and \((P', p', q')\) of \(\mathcal{T}_{n}(M, N)\) are said to be linked if there exists a morphism \(u: P \to P'\) such that \(p = p' \circ u\) and \(q' = q \circ 'u\) ; we denote by \(T_{n}(M, N)\) the quotient of \(\mathcal{T}_{n}(M, N)\) by the finest equivalence relation for which two related elements are equivalent (cf.~II, p.~52, exercise 9). a) Show that \(T_{n}(M, N)\) is identified with Tor \(_{n}^{A}\) \((M, N)\) .
\end{enumerate}

\begin{enumerate}
\def\labelenumi{\alph{enumi})}
\setcounter{enumi}{1}
\tightlist
\item
  Let \((\mathcal{E}):0\to \mathrm{M}^{\prime}\stackrel {i}{\to}\mathrm{M}\stackrel {\pi}{\to}\mathrm{M}^{\prime \prime}\to 0\) be an exact sequence of right A-modules. Describe the connecting homomorphism \(\partial :\mathbf{T}_n(\mathbf{M}'',\mathbf{N})\to \mathbf{T}_{n - 1}(\mathbf{M}',\mathbf{N})\) deduced from \(\partial (\mathcal{E},\mathbf{N})\) using the previous identifications.
\end{enumerate}

\((\mathrm{If}(\mathbf{P},p,q)\in\mathcal{T}_{n}(\mathbf{M}^{\prime\prime},\mathbf{N}),\text{show that we have}\partial(\mathbf{P},p,q)=(\overline{\mathbf{P}},\overline{p},\overline{q}),\text{where}\overline{\mathbf{P}}\text{ is the complex obtained by replacing } \mathbf{P}_{n}\text{ by 0 in } \mathbf{P},\overline{p}\text{ is the morphism deduced from } p\text{ and }(\overline{q})^{n-1}: \mathbf{P}_{n-1}^{*}\to\mathbf{M}^{\prime}\text{ is equal to }u^{n-1},u\text{ being a morphism from } \mathbf{P}^{*}\text{ into the complex (Con (i)) (n) such that }\pi\circ u^{n}=q^{n}.)\)

¶ 7) Let A' be an (associative and unital) k-algebra and B = A ⊗k A'.

\begin{enumerate}
\def\labelenumi{\alph{enumi})}
\tightlist
\item
  Let M be a left A-module, Q a right A-module, M' a left A'-module, Q' a right A'-module. Let (P, p) (resp. (P', p'), resp. (R, r)) be a projective resolution of the A-module M (resp. of the A'-module M', resp. of the B-module M ⊗\_k M'). Show that there exists a morphism of B-complexes
\end{enumerate}

\[
u: \mathbf {P} \otimes_ {k} \mathbf {P} ^ {\prime} \rightarrow \mathbf {R},
\]

unique up to homotopy, such that \(r \circ u = p \otimes p'\) . Define, using \(u\) , a graded \(k\) -homomorphism of degree zero

\[
\top : \operatorname{Tor} ^ {\mathbf {A}} (\mathrm{Q}, \mathrm{M}) \otimes_ {k} \operatorname{Tor} ^ {\mathbf {A} ^ {\prime}} (\mathrm{Q} ^ {\prime}, \mathrm{M} ^ {\prime}) \rightarrow \operatorname{Tor} ^ {\mathbf {B}} (\mathrm{Q} \otimes_ {k} \mathrm{Q} ^ {\prime}, \mathrm{M} \otimes_ {k} \mathrm{M} ^ {\prime}).
\]

Show that \(\top\) does not depend on the choice of resolutions P, P', R and of the morphism \(u\) . If \(a \in \operatorname{Tor}^{\mathbf{A}}(\mathbf{Q}, \mathbf{M})\) and \(a' \in \operatorname{Tor}^{\mathbf{A}'}(\mathbf{Q}', \mathbf{M}')\) , we set \(\top(a \otimes a') = a \top a'\) .

\begin{enumerate}
\def\labelenumi{\alph{enumi})}
\setcounter{enumi}{1}
\tightlist
\item
  We assume moreover that the \(k\) -modules \(\mathbf{A}\) and \(\mathbf{A}'\) are projective and that one has \(\operatorname{Tor}_n^k (\mathbf{M},\mathbf{M}') = 0\) for \(n > 0\) . Show that \(\mathbf{P} \otimes_k \mathbf{P}'\) is then a projective resolution of the \(\mathbf{B}\) -module \(\mathbf{M} \otimes_k \mathbf{M}'\) ; deduce from this, for every \(\mathbf{A}\) -module \(\mathbf{N}\) and for every \(\mathbf{A}'\) -module \(\mathbf{N}'\) , a graded \(k\) -homomorphism of degree zero
\end{enumerate}

\[
\vee : \operatorname{Ext} _ {\mathbf {A}} (\mathbf {M}, \mathbf {N}) \otimes_ {k} \operatorname{Ext} _ {\mathbf {A} ^ {\prime}} (\mathbf {M} ^ {\prime}, \mathbf {N} ^ {\prime}) \rightarrow \operatorname{Ext} _ {\mathbf {B}} (\mathbf {M} \otimes_ {k} \mathbf {M} ^ {\prime}, \mathbf {N} \otimes_ {k} \mathbf {N} ^ {\prime}).
\]

Show that \(\vee\) is independent of the choice of resolutions \(\mathbf{P}\) and \(\mathbf{P}'\) . If \(b \in \operatorname{Ext}_{\mathbf{A}}(\mathbf{M}, \mathbf{N})\) and \(b' \in \operatorname{Ext}_{\mathbf{A}'}(\mathbf{M}', \mathbf{N}')\) , we write \(\vee (b \otimes b') = b \vee b'\) .

\begin{enumerate}
\def\labelenumi{\alph{enumi})}
\setcounter{enumi}{2}
\tightlist
\item
  Let A'' be a third k-algebra (associative and unital), M'' and N'' left A''-modules, Q'' a right A''-module. We identify the k-algebras A ⊗\_k (A' ⊗\_k A'') and (A ⊗\_k A') ⊗\_k A'', as well as the k-modules M ⊗\_k (M' ⊗\_k M'') and (M ⊗\_k M') ⊗\_k M'', N ⊗\_k (N' ⊗\_k N'') and (N ⊗\_k N') ⊗\_k N'', Q ⊗\_k (Q' ⊗\_k Q'') and (Q ⊗\_k Q') ⊗\_k Q''. Prove the formulas
\end{enumerate}

\[
\begin{array}{l l} a \top (a ^ {\prime} \top a ^ {\prime \prime}) = (a \top a ^ {\prime}) \top a ^ {\prime \prime} & \text {for} \quad a \in \operatorname{Tor} ^ {\mathbf {A}} (\mathrm{Q}, \mathrm{M}), \quad a ^ {\prime} \in \operatorname{Tor} ^ {\mathbf {A} ^ {\prime}} (\mathrm{Q} ^ {\prime}, \mathrm{M} ^ {\prime}), \\ & \qquad \qquad \qquad \qquad \qquad \qquad \qquad \qquad \qquad \qquad \qquad \qquad \qquad \qquad a ^ {\prime \prime} \in \operatorname{Tor} ^ {\mathbf {A} ^ {\prime \prime}} (\mathrm{Q} ^ {\prime \prime}, \mathrm{M} ^ {\prime \prime}) \end{array}
\]

\[
\begin{array}{r l} b \vee (b ^ {\prime} \vee b ^ {\prime \prime}) = (b \vee b ^ {\prime}) \vee b ^ {\prime \prime} & \text {for} \quad b \in \operatorname{Ext} _ {\mathbf {A}} (\mathbf {M}, \mathbf {N}), \quad b ^ {\prime} \in \operatorname{Ext} _ {\mathbf {A}} (\mathbf {M} ^ {\prime}, \mathbf {N} ^ {\prime}), \\ & \qquad \qquad \qquad \qquad \qquad \qquad b ^ {\prime \prime} \in \operatorname{Ext} _ {\mathbf {A} ^ {\prime \prime}} (\mathbf {M} ^ {\prime \prime}, \mathbf {N} ^ {\prime \prime}). \end{array}
\]

\begin{enumerate}
\def\labelenumi{\alph{enumi})}
\setcounter{enumi}{3}
\tightlist
\item
  Let \(\sigma : \mathrm{Tor}^{\mathbf{A} \otimes_{\mathbf{k}} \mathbf{A}'}(\mathbf{Q} \otimes_{\mathbf{k}} \mathbf{Q}', \mathbf{M} \otimes_{\mathbf{k}} \mathbf{M}') \to \mathrm{Tor}^{\mathbf{A}' \otimes_{\mathbf{k}} \mathbf{A}}(\mathbf{Q}' \otimes_{\mathbf{k}} \mathbf{Q}, \mathbf{M}' \otimes_{\mathbf{k}} \mathbf{M})\) and
\end{enumerate}

\[
\tau : \operatorname{Ext} _ {\mathbf {A} \otimes_ {k} \mathbf {A} ^ {\prime}} (\mathrm{M} \otimes_ {k} \mathrm{M} ^ {\prime}, \mathrm{N} \otimes_ {k} \mathrm{N} ^ {\prime}) \rightarrow \operatorname{Ext} _ {\mathbf {A} ^ {\prime} \otimes_ {k} \mathbf {A}} (\mathrm{M} ^ {\prime} \otimes_ {k} \mathrm{M}, \mathrm{N} ^ {\prime} \otimes_ {k} \mathrm{N})
\]

the isomorphisms deduced from the commutation isomorphisms

\[
\mathrm{A} \otimes_ {k} \mathrm{A} ^ {\prime} \rightarrow \mathrm{A} ^ {\prime} \otimes_ {k} \mathrm{A}, \quad \mathrm{M} \otimes_ {k} \mathrm{M} ^ {\prime} \rightarrow \mathrm{M} ^ {\prime} \otimes_ {k} \mathrm{M}, \quad \mathrm{N} \otimes_ {k} \mathrm{N} ^ {\prime} \rightarrow \mathrm{N} ^ {\prime} \otimes_ {k} \mathrm{N}, \quad \mathrm{Q} \otimes_ {k} \mathrm{Q} ^ {\prime} \rightarrow \mathrm{Q} ^ {\prime} \otimes_ {k} \mathrm{Q}.
\]

Prove the formulas.

\[
\begin{array}{l l} \sigma (a \top a ^ {\prime}) = (- 1) ^ {p q} a ^ {\prime} \top a & \text { for } \quad a \in \operatorname{Tor} _ {p} ^ {A} (Q, M), \quad a ^ {\prime} \in \operatorname{Tor} _ {q} ^ {A ^ {\prime}} (Q ^ {\prime}, M ^ {\prime}) \\ \tau (b \vee b ^ {\prime}) = (- 1) ^ {p q} b ^ {\prime} \vee b & \text { for } \quad b \in \operatorname{Ext} _ {A} ^ {p} (M, N), \quad b ^ {\prime} \in \operatorname{Ext} _ {A ^ {\prime}} ^ {q} (M ^ {\prime}, N ^ {\prime}). \end{array}
\]

\begin{enumerate}
\def\labelenumi{\alph{enumi})}
\setcounter{enumi}{4}
\tightlist
\item
  Let \((\mathcal{E}):0\to \mathbf{M}_2\to \mathbf{M}_1\to \mathbf{M}_0\to 0\) be an exact sequence of A-modules, such that the sequence
\end{enumerate}

\[
(\mathcal {E} \otimes \mathrm{M} ^ {\prime}): 0 \rightarrow \mathrm{M} _ {2} \otimes_ {k} \mathrm{M} ^ {\prime} \rightarrow \mathrm{M} _ {1} \otimes_ {k} \mathrm{M} ^ {\prime} \rightarrow \mathrm{M} _ {0} \otimes_ {k} \mathrm{M} ^ {\prime} \rightarrow 0
\]

be exact. Let δ (resp. Δ) denote the connecting homomorphism δ(Q, ℰ) (resp. δ(Q ⊗ₖ Q′, ℰ ⊗ M′)). Show that one has (δa) ⊤ a′ = Δ(a ⊗ a′) for a ∈ Tor\^{}A(Q, M₀), a′ ∈ Tor\^{}A(Q', M'). Prove the analogous formulas for the product ∨.

\begin{enumerate}
\def\labelenumi{\alph{enumi})}
\setcounter{enumi}{5}
\tightlist
\item
  If k is a field, show that \(\top\) is an isomorphism of graded k-modules; the same holds for ∨ if, moreover, A and A' are left noetherian and the A-module M and the A'-module M' are finitely generated. g) If one identifies the modules \(\operatorname{Tor}_{n}^{\mathbf{A}}(\mathbf{Q},\mathbf{M})\) with the sets \(\operatorname{T}_{n}(\mathbf{Q},\mathbf{M})\) defined in Exercise 6, prove the formula \((\mathbf{P},p,q)\top(\mathbf{P}',p',q')=(\mathbf{P}\otimes_{k}\mathbf{P}',p\otimes p',q\otimes q')\) .
\end{enumerate}

\begin{enumerate}
\def\labelenumi{\arabic{enumi})}
\setcounter{enumi}{7}
\tightlist
\item
  We keep the notation of Exercise 7; we assume that the \(k\) -modules A and A' are projective.
\end{enumerate}

\begin{enumerate}
\def\labelenumi{\alph{enumi})}
\item
  We assume that one has \(\mathrm{Tor}_{i}^{k}(\mathbf{M},\mathbf{M}^{\prime}) = 0\) for \(i > 0\) . Show that the map \(b\mapsto b\vee 1_{\mathbf{M}'}\) defines a graded \(k\) -homomorphism of degree zero from \(\operatorname{Ext}_{\mathbf{A}}(\mathbf{M},\mathbf{N})\) to \(\operatorname{Ext}_{\mathbf{B}}(\mathbf{M}\otimes_{k}\mathbf{M}',\mathbf{N}\otimes_{k}\mathbf{M}')\) , which induces in degree zero the canonical homomorphism \(b\mapsto b\otimes 1_{\mathbf{M}'}\) .
\item
  We assume that the k-module M' is flat. Let
\end{enumerate}

\[
0 \rightarrow \mathrm{N} \rightarrow \mathrm{R} _ {n} \rightarrow \dots \rightarrow \mathrm{R} _ {1} \rightarrow \mathrm{M} \rightarrow 0
\]

be an exact sequence of A-modules, of class \(b \in \operatorname{Ext}_{\mathbf{A}}^{n}(\mathbf{M}, \mathbf{N})\) . Show that the exact sequence of B-modules

\[
0 \rightarrow \mathrm{N} \otimes_ {k} \mathrm{M} ^ {\prime} \rightarrow \mathrm{R} _ {n} \otimes_ {k} \mathrm{M} ^ {\prime} \rightarrow \dots \rightarrow \mathrm{R} _ {1} \otimes_ {k} \mathrm{M} ^ {\prime} \rightarrow \mathrm{M} \otimes_ {k} \mathrm{M} ^ {\prime} \rightarrow 0
\]

has class \(b \vee 1_{\mathbf{M}'}\) in \(\operatorname{Ext}_{\mathbf{B}}^{n}(\mathbf{M} \otimes_{k} \mathbf{M}', \mathbf{N} \otimes_{k} \mathbf{N}')\) .

\begin{enumerate}
\def\labelenumi{\alph{enumi})}
\setcounter{enumi}{2}
\tightlist
\item
  We assume that the k-modules M, M', N, N' are flat. Show that we have:
\end{enumerate}

\[
b \vee b ^ {\prime} = (b \vee 1 _ {N ^ {\prime}}) \circ (1 _ {M} \vee b ^ {\prime}) = (- 1) ^ {p q} (1 _ {N} \vee b ^ {\prime}) \circ (b \vee 1 _ {M ^ {\prime}})
\]

for \(b \in \operatorname{Ext}_{\mathbf{A}}^{p}(\mathbf{M}, \mathbf{N})\) , \(b' \in \operatorname{Ext}_{\mathbf{A}'}^{q}(\mathbf{M}', \mathbf{N}')\) .

\begin{enumerate}
\def\labelenumi{\arabic{enumi})}
\setcounter{enumi}{8}
\tightlist
\item
  We assume the ring A is commutative. Let B, C be two associative and unital A-algebras; we denote by \(m_{\mathbf{B}}: \mathbf{B} \otimes_{\mathbf{A}} \mathbf{B} \to \mathbf{B}\) the multiplication of B, and \(m_{\mathbf{C}}\) the multiplication of C.
\end{enumerate}

\begin{enumerate}
\def\labelenumi{\alph{enumi})}
\tightlist
\item
  Show that the homomorphism
\end{enumerate}

\[
\operatorname{Tor} \left(m _ {\mathrm{B}}, m _ {\mathrm{C}}\right) \circ \top : \operatorname{Tor} ^ {\mathrm{A}} (\mathrm{B}, \mathrm{C}) \otimes_ {\mathrm{A}} \operatorname{Tor} ^ {\mathrm{A}} (\mathrm{B}, \mathrm{C}) \rightarrow \operatorname{Tor} ^ {\mathrm{A}} (\mathrm{B}, \mathrm{C})
\]

(cf.~exercise 7) defines on Tor \(^{A}\) (B, C) a structure of graded A-algebra, associative and unital. b) Let S be a differential graded A-algebra and \(p : S_{0} \to B\) a homomorphism of A-algebras such that (S, p) is a free resolution of B (cf. X, p. 183, exercise 18). Show that the isomorphism

\[
\psi (\mathbf {S}, \mathbf {C}): \operatorname{Tor} ^ {\mathbf {A}} (\mathbf {B}, \mathbf {C}) \rightarrow \mathbf {H} (\mathbf {S} \otimes_ {\mathbf {A}} \mathbf {C})
\]

is an isomorphism of A-algebras if one endows H(S ⊗A C) with the algebra structure deduced from the differential graded algebra structure of S ⊗A C (X, p.~183, exercise 18, b) and c)).

\begin{enumerate}
\def\labelenumi{\alph{enumi})}
\setcounter{enumi}{2}
\tightlist
\item
  If the A-algebras B and C are commutative, prove that the graded algebra Tor \(^{A}\) (B, C) is alternating (cf. X, p. 183, exercise 19).
\end{enumerate}

\begin{enumerate}
\def\labelenumi{\arabic{enumi})}
\setcounter{enumi}{9}
\tightlist
\item
  Let B be a k-bialgebra (III, p.~148), which we regard as an augmented k-algebra (X, p.~196, exercise 27) via the counit \(\beta: \mathbf{B} \to k\) ; we assume that the \(k\) -module B is projective.
\end{enumerate}

\begin{enumerate}
\def\labelenumi{\alph{enumi})}
\tightlist
\item
  Let M, N be two left B-modules. The homomorphism ∨ of Exercise 7 is identified with a graded homomorphism of degree zero
\end{enumerate}

\[
\mathrm{H} ^ {\bullet} (\mathbf {B}, \gamma ; \mathbf {M}) \otimes_ {k} \mathrm{H} ^ {\bullet} (\mathbf {B}, \gamma ; \mathbf {N}) \rightarrow \mathrm{H} ^ {\bullet} (\mathbf {B} \otimes_ {k} \mathbf {B}, \gamma \otimes \gamma ; \mathbf {M} \otimes_ {k} \mathbf {N});
\]

by composing with the homomorphism induced by the comultiplication, deduce from this a graded homomorphism of degree zero \(\cup: H'(B, \gamma; M) \otimes_k H'(B, \gamma; N) \to H'(B, \gamma; M \otimes_k N)\) .

\begin{enumerate}
\def\labelenumi{\alph{enumi})}
\setcounter{enumi}{1}
\item
  Let C be an associative unital k-algebra, which we equip with the B-module structure induced by γ. Show that the homomorphism ∪ defines on H'(B, γ ; C) a structure of associative unital graded k-algebra; if C is commutative and the bialgebra B is cocommutative, the algebra H'(B, γ ; C) is anticommutative.
\item
  We take C = k. Show that the product ∪ coincides with the composition product on \(\operatorname{Ext}_{\mathbf{B}}(k, k)\) . d) Let G be a group, M and N two \(\mathbf{Z}^{(\mathbf{G})}\) -modules. We endow M ⊗\_z N with the structure of \(\mathbf{Z}^{(\mathbf{G})}\) -module defined by \(g(x \otimes y) = gx \otimes gy\) for \(g \in G\) , \(x \in M\) , \(y \in N\) . Show that the map
\end{enumerate}

\[
\cup : \mathrm{H} ^ {\bullet} (\mathrm{G}, \mathrm{M}) \otimes_ {\mathbf {z}} \mathrm{H} ^ {\bullet} (\mathrm{G}, \mathrm{N}) \rightarrow \mathrm{H} ^ {\bullet} (\mathrm{G}, \mathrm{M} \otimes_ {\mathbf {z}} \mathrm{N})
\]

is the homomorphism induced on homology by the graded homomorphism of degree zero

\[
u: \mathrm{C} ^ {*} (\mathrm{G}, \mathrm{M}) \otimes_ {\mathbf {z}} \mathrm{C} ^ {*} (\mathrm{G}, \mathrm{N}) \rightarrow \mathrm{C} ^ {*} (\mathrm{G}, \mathrm{M} \otimes_ {\mathbf {z}} \mathrm{N})
\]

defined by \(u(f \otimes h)(g_1, \ldots, g_{p+q}) = f(g_1, \ldots, g_p) \otimes g_1 \ldots g_p h(g_{p+1}, \ldots, g_{p+q})\) for

\[
f \in \mathbf {C} ^ {p} (\mathbf {G}, \mathbf {M}), \quad h \in \mathbf {C} ^ {q} (\mathbf {G}, \mathbf {N}), \quad g _ {1} \dots g _ {p + q} \in \mathbf {G}.
\]

(Use Exercise 5, p.~185.)

¶ 11) Suppose that the bigebra B admits an inversion i (cf.~III, p.~198, exercise 4). Let M and N be B-modules. The k-module M ⊗k N (resp. Homk(M, N)) has a natural structure of B ⊗k B-module (resp. B ⊗k B°-module); one endows it with the B-module structure deduced from the homomorphism c (resp. (1B ⊗ i) ∘ c).

\begin{enumerate}
\def\labelenumi{\alph{enumi})}
\tightlist
\item
  Show that the homomorphism u (Exercise 10) defines a graded homomorphism of degree zero
\end{enumerate}

\[
\cup : \mathrm{H} ^ {*} (\mathbf {B}, \gamma ; \operatorname{Hom} _ {k} (\mathbf {M}, \mathbf {N})) \otimes_ {k} \mathrm{H} ^ {*} (\mathbf {B}, \gamma ; \mathbf {M}) \rightarrow \mathrm{H} ^ {*} (\mathbf {B}, \gamma ; \mathbf {N}).
\]

\begin{enumerate}
\def\labelenumi{\alph{enumi})}
\setcounter{enumi}{1}
\tightlist
\item
  Let \(0 \to M \to R_n \to \cdots \to R_1 \to k \to 0\) be an exact sequence of B-modules, with \(R_i\) projective ( \(1 \leqslant i \leqslant n\) ); let \(\theta\) be its class in \(\operatorname{Ext}_B^n(k, M) = H^n(B, \gamma; M)\) . Show that for \(p \geqslant 1\) , the map \(x \mapsto x \cup \theta\) is a \(k\) -isomorphism of \(H^p(B, \gamma; \operatorname{Hom}_k(M, N))\) onto \(H^{p+n}(B, \gamma; N)\) .
\end{enumerate}

(Show that this map factors through

\[
\mathrm{H} ^ {p} (\mathrm{B}, \gamma ; \operatorname{Hom} _ {k} (\mathrm{M}, \mathrm{N})) \rightarrow \operatorname{Ext} _ {\mathrm{B}} ^ {p} (\mathrm{M}, \mathrm{N}) \xrightarrow {\theta} \operatorname{Ext} _ {\mathrm{B}} ^ {p + n} (k, \mathrm{N}).
\]

\begin{enumerate}
\def\labelenumi{\alph{enumi})}
\setcounter{enumi}{2}
\tightlist
\item
  Let G be a group, N a \(\mathbf{Z}^{(G)}\) -module. Let F be a free group, \(p: F \to G\) a surjective homomorphism, R = Ker( \(p\) ), \(\overline{R} = R/(R, R)\) , \(\overline{F} = F/(R, R)\) . The extension \(\overline{R} \to \overline{F} \to G\) defines a class \(\varepsilon \in H^2(G, \overline{R})\) ; prove that the map \(x \mapsto x \cup \varepsilon\) from \(H^p(G, \operatorname{Hom}_{\mathbf{Z}}(\overline{R}, N))\) to \(H^{p+2}(G, N)\) is an isomorphism for \(p \geqslant 1\) (use exercise 3, d), p.~179).
\end{enumerate}

